% Part: sets-functions-relations
% Chapter: functions
% Section: functions-relations

\documentclass[../../../include/open-logic-section]{subfiles}


\begin{document}

\olfileid{sfr}{fun}{rel}

\olsection{সম্পর্ক হিসেবে অপেক্ষক}

\begin{explain}
যে অপেক্ষক~$A$-এর !!{element}s-কে~$B$-এর !!{element}s-এ পাঠায়, সেটি স্পষ্টতই $A$ ও~$B$-এর মধ্যে একটি সম্পর্ক নির্ধারণ করে: $x$ ও $y$-এর মধ্যে সম্পর্কটি থাকবে যদি এবং কেবল যদি $f(x) = y$ হয়। বস্তুত, গণিতের মৌলিক নির্মাণ-উপকরণের সংখ্যা কমাতে আগ্রহী হলে, যেমন, আমরা অপেক্ষক~$f$-কে এই সম্পর্কের সঙ্গে, অর্থাৎ ক্রমযুগলের একটি সেটের সঙ্গে \emph{অভিন্ন} বলে ধরতেও পারি। তখন প্রশ্ন ওঠে: কোন কোন সম্পর্ক এভাবে অপেক্ষক নির্ধারণ করে?
\end{explain}

\begin{defn}[অপেক্ষকের গ্রাফ] ধরি, $f\colon A \to B$ একটি অপেক্ষক।
নিচের সংজ্ঞা দ্বারা নির্ধারিত সম্পর্ক $R_f \subseteq A \times B$-কে~$f$-এর \emph{গ্রাফ} বলা হয়:
\[
R_f = \Setabs{\tuple{x,y}}{f(x) = y}.
\]
\end{defn}

\begin{explain}
সদস্যভিত্তিক সমতার নীতি অনুসারে একটি অপেক্ষকের গ্রাফ অনন্যভাবে নির্ধারিত। তদুপরি, সেটের সদস্যভিত্তিক সমতার নীতি থেকেই অপেক্ষকের মানভিত্তিক সমতার অন্তর্নিহিত নীতিটি সরাসরি সমর্থিত হয়: $f$ ও~$g$-এর সংজ্ঞাক্ষেত্র এবং সহসংজ্ঞাক্ষেত্র একই হলে, সব ইনপুটে তাদের মান এক হলে তারা অভিন্ন।

একইভাবে, কোনো সম্পর্ক ``অপেক্ষকধর্মী'' হলে সেটি একটি অপেক্ষকের গ্রাফ।
\end{explain}

\begin{prop}\ollabel{prop:graph-function}
ধরি, $R \subseteq A \times B$-এর জন্য নিচের শর্তগুলি পূরণ হয়:
\begin{enumerate}
\item $Rxy$ ও $Rxz$ হলে $y = z$; এবং
\item প্রত্যেক $x \in A$-এর জন্য এমন কোনো $y \in B$ আছে, যাতে $\tuple{x,
y} \in R$।
\end{enumerate}
তাহলে $R$ হলো সেই অপেক্ষক $f\colon A \to B$-এর গ্রাফ, যার সংজ্ঞা: $f(x) = y$ যদি এবং কেবল যদি $Rxy$।
\end{prop}

\begin{proof}
ধরি, এমন একটি $y$ আছে, যাতে $Rxy$। যদি অন্য কোনো $z \neq
y$-এর জন্য $Rxz$ হতো, তবে~$R$-এর উপর আরোপিত শর্ত লঙ্ঘিত হতো। সুতরাং $Rxy$ হয় এমন কোনো $y$ থাকলে, সেই $y$ অনন্য; তাই $f$ সুসংজ্ঞায়িত। স্পষ্টতই $R_f = R$।
\end{proof}

\begin{explain}
প্রত্যেক অপেক্ষক $f\colon A \to B$-এর একটি গ্রাফ আছে, অর্থাৎ $f(x) = y$ দ্বারা নির্ধারিত, $A \times B$-এর অন্তর্ভুক্ত A ও B-এর মধ্যে একটি সম্পর্ক আছে। অন্যদিকে, \olref{prop:graph-function}-এ দেওয়া ধর্মগুলি আছে এমন প্রত্যেক সম্পর্ক~$R
\subseteq A \times B$ কোনো অপেক্ষক~$f \colon A \to
B$-এর গ্রাফ। অপেক্ষক ও তাদের গ্রাফের এই ঘনিষ্ঠ যোগসূত্রের ফলে আমরা একটি অপেক্ষককে কেবল তার গ্রাফ হিসেবেই ভাবতে পারি। অন্যভাবে বললে, অপেক্ষককে বিশেষ কিছু সম্পর্কের সঙ্গে, অর্থাৎ টিউপলের বিশেষ কিছু সেটের সঙ্গে অভিন্ন বলে ধরা যায়। \oliflabeldef{sfr:rel:ref:sec}{তবে লক্ষ করো, এই ``অভিন্ন বলে ধরা'' কথাটির তাৎপর্য \olref[sfr][rel][ref]{sec}-এর মতোই: এটি অপেক্ষকের অধিবিদ্যা সম্পর্কে কোনো দাবি নয়, বরং অপেক্ষককে বিশেষ কিছু সেট হিসেবে \emph{বিবেচনা করা} সুবিধাজনক—এই পর্যবেক্ষণ। এর একটি সুবিধা হলো, }{}এখন আমরা সম্পর্কের উপর যেসব ক্রিয়া করেছিলাম, অপেক্ষকের উপরও অনুরূপ ক্রিয়া করার কথা বিবেচনা করতে পারি (দেখো \olref[sfr][rel][ops]{sec})। বিশেষ করে: % OLFUN-004 TARGET-CORRECTION-BEGIN
\par\smallskip\noindent\textnormal{\small\textbf{উৎস-সংশোধন (OLFUN-004):} ইংরেজি উৎসে এক স্থানে গ্রাফকে ``$A\times B$-এর উপর সম্পর্ক'' বলা হয়েছে। গ্রন্থটির নিজস্ব সংজ্ঞা অনুযায়ী সেটি আক্ষরিকভাবে $(A\times B)^2$-এর উপসেট বোঝাত; সঠিক অর্থ হলো $A$ ও~$B$-এর মধ্যে সম্পর্ক, অর্থাৎ $A\times B$-এর উপসেট।} % OLFUN-004 TARGET-CORRECTION-END
\end{explain}

\begin{defn}\ollabel{defn:funimage}
ধরি, $f \colon A \to B$ একটি অপেক্ষক এবং $C\subseteq A$।

প্রত্যেক $x \in C$-এর জন্য $(\funrestrictionto{f}{C})(x) = f(x)$ দ্বারা সংজ্ঞায়িত অপেক্ষক~$\funrestrictionto{f}{C}\colon C \to B$-কে~$C$-তে~$f$-এর \emph{সীমাবদ্ধন} বলা হয়। অন্যভাবে বললে, $\funrestrictionto{f}{C} = \Setabs{\tuple{x, y} \in R_f}{x \in
C}$।

সেট~$C$-তে~$f$-এর \emph{প্রয়োগ} হলো $\funimage{f}{C} =
\Setabs{f(x)}{x \in C}$। একে~$f$-এর অধীনে~$C$-এর \emph{প্রতিবিম্ব}ও বলা হয়।
\end{defn}

\begin{explain}
এই সংজ্ঞাগুলি থেকে পাওয়া যায় যে যেকোনো অপেক্ষক~$f$-এর জন্য $\ran{f} =
\funimage{f}{\dom{f}}$।
\oliflabeldef{sfr:rel:ops:sec}{\olref[sfr][rel][ops]{sec}-এর ধারণাগুলির সঙ্গে এগুলি ঘনিষ্ঠভাবে সম্পর্কিত; তবে অপেক্ষকের সীমাবদ্ধনে শুধু ইনপুট স্থানাঙ্কটি C-তে সীমিত হয়, যেখানে আগের সম্পর্কের সীমাবদ্ধনে উভয় স্থানাঙ্কই C-তে সীমিত হয়। আরও দুটি ক্রিয়া—বিপরীত ও আপেক্ষিক গুণফল—কিছুটা বিস্তৃত আলোচনা দাবি করে। \olref[inv]{sec} ও \olref[cmp]{sec}-এ সেই আলোচনা করব।}{} % OLFUN-005 TARGET-CORRECTION-BEGIN
\par\smallskip\noindent\textnormal{\small\textbf{উৎস-সংশোধন (OLFUN-005):} ইংরেজি উৎসের ``exactly as one would expect'' কথাটি যোগ্যতাসহ বলা হয়েছে। এখানে অপেক্ষকের গ্রাফ $\Setabs{\tuple{x,y}\in R_f}{x\in C}$; আগের সম্পর্কের সীমাবদ্ধন $R_f\cap C^2$। ফলে দুটি সংজ্ঞা সাধারণভাবে এক নয়।} % OLFUN-005 TARGET-CORRECTION-END
\end{explain}

\end{document}
